\documentclass[10pt,a4paper]{amsart}
\usepackage[margin=19mm]{geometry}
\usepackage{amsmath,amssymb,mathtools}
\usepackage[scr=rsfs,cal=euler]{mathalfa}
\usepackage{xfrac}
\usepackage[unicode,hidelinks,bookmarks=false]{hyperref}
\newcommand{\ie}{i.e.\ }
\newcommand{\cf}{cf.\ }
\newcommand{\resp}{respectively\ }
\newcommand{\Subsection}{\subsection}
\providecommand{\nobreakdash}{\nobreak\hbox{-}}
\setlength{\emergencystretch}{2em}
\multlinegap0pt
\usepackage{amssymb}
\newtheorem{lemma}{Lemma}
\newtheorem{proposition}{Proposition}
\newcommand{\IQbar}{{\overline{\QQ}}}
\def\QQ{\mathbb{Q}}
\newcommand{\ssm}{\setminus}
\renewcommand{\subset}{\subseteq} %%% Some people think \subset
\title{The relative Manin--Mumford conjecture}
\author{Ziyang Gao, Philipp Habegger}
\date{Source: arXiv:2303.05045v3 (CC BY 4.0)}
\begin{document}
\maketitle

\noindent\emph{Source and licence.} The relative Manin--Mumford conjecture. Version arXiv:2303.05045v3 (CC BY 4.0), distributed by arXiv under the Creative Commons Attribution 4.0 International licence, http://creativecommons.org/licenses/by/4.0/, as recorded in the arXiv licence metadata for this exact version and shown on the abstract page https://arxiv.org/abs/2303.05045v3. Full licence notice: \texttt{licenses/arxiv-2303.05045-CC-BY-4.0.txt}.\par\smallskip

\noindent\textit{Editorial scope.} Counted unit: the article's proposition PropUML (source0.tex lines 1589--1593). The article's complete original proof of it is delivered with it (source0.tex lines 1594--1604, one \texttt{proof} environment of 11 lines). No mathematics is written, repaired or re-derived: the statement and the proof are the article's own lines, in the article's own notation, and no retained line is substituted or re-flowed.  Retained uncounted as the source-local context the proof needs: lines 1483--1495.  Nothing was omitted from any retained span, so the package carries no omission record.  No retained line contains a citation, so the wrapper carries no bibliography and no bibliography entry.  No retained line contains a figure environment, an image-inclusion command or drawing code, so no image is delivered and the package declares no asset dependency.  Editorial formatting only, in addition: an AMS \texttt{amsart} preamble that restates the article's own declarations and macros used by the retained lines, namely the environments \texttt{lemma}, \texttt{proposition} on separate counters, together with the standard packages those lines need.  The retained environments print in this document as Lemma 1, Proposition 1 in order of appearance, whereas the article prints the same items with the numbering of its own full document.  The complete native source of the article as distributed in the arXiv e-print for this exact version is bundled unchanged as \texttt{sources/138395/source0.tex}.
\begin{lemma}\label{LemUML}
  Let $S$ be a regular, irreducible, locally closed subvariety of
  $\mathbb{M}_g$. Then there exists a constant $c = c(S) > 0$ such that
  for all $s\in S(\IQbar)$ there is $\Xi_s\subset
  \mathfrak{C}_s(\IQbar)$ with $\#\Xi_s < c$ and with the following
  property.
  For all $x\in \mathfrak{C}_s(\IQbar)\ssm \Xi_s$ we have
  % \begin{enumerate}
  % \item[(i)] Either $x$ belongs to a subset of $\mathfrak{C}_s(\IQbar)$ of cardinality at most $c$,
  % \item[(ii)] or
    $\#\{y \in \mathfrak{C}_s(\IQbar) : y-x \in (\mathfrak{J}_s)_{\mathrm{tor}} \} < c$.
%  \end{enumerate}
\end{lemma}

\begin{proposition}\label{PropUML}
  Let $S$ and $c=c(S)$ be as in Lemma~\ref{LemUML}. For
  all $s\in S(\IQbar)$ and all $x \in \mathfrak{C}_s(\IQbar)$ we
  have  $\#\{y \in \mathfrak{C}_s(\IQbar) : y-x \in (\mathfrak{J}_s)_{\mathrm{tor}} \} < c$.
\end{proposition}

\begin{proof}
  Let $s\in S(\IQbar)$ and $x\in \mathfrak{C}_s(\IQbar)$.
  Suppose $y_1,\ldots,y_n\in\mathfrak{C}_s(\IQbar)$ are pairwise
  distinct with $y_i-x$ torsion for all $i$.
  If $n\ge c$, then some $y_i$, say $y_1$, lies in
  $\mathfrak{C}_s(\IQbar)\ssm \Xi_s$ with $\Xi_s$ the set from
  Lemma~\ref{LemUML}.
  Now $y_i-y_1 = (y_i-x)-(y_1-x)$ are $n$  torsion points
  of $\mathfrak{J}_s(\IQbar)$ for $i\in\{1,\ldots,n\}$. So
  Lemma~\ref{LemUML} implies $n<c$, a contradiction. We conclude $n<c$. 
\end{proof}

\end{document}
